% Szablon Pracy Dyplomowej w LaTeX (Standard 2027)
% Kompatybilny z Tectonic, Overleaf oraz silnikiem pdfLaTeX / XeLaTeX

\documentclass[12pt,a4paper,twoside]{report}
\usepackage[utf8]{inputenc}
\usepackage[T1]{fontenc}
\usepackage{lmodern}
\usepackage{amsmath,amssymb}
\usepackage{graphicx}
\usepackage{listings}
\usepackage{xcolor}
\usepackage{geometry}
\usepackage{hyperref}
\usepackage[style=ieee,backend=biber]{biblatex}

\geometry{
  a4paper,
  left=30mm,
  right=20mm,
  top=25mm,
  bottom=25mm
}

\definecolor{codegray}{rgb}{0.5,0.5,0.5}
\definecolor{codepurple}{rgb}{0.58,0,0.82}
\definecolor{backcolour}{rgb}{0.96,0.96,0.96}

\lstdefinestyle{mystyle}{
    backgroundcolor=\color{backcolour},   
    commentstyle=\color{codegray},
    keywordstyle=\color{blue},
    numberstyle=\tiny\color{codegray},
    stringstyle=\color{codepurple},
    basicstyle=\ttfamily\footnotesize,
    breakatwhitespace=false,         
    breaklines=true,                 
    captionpos=b,                    
    keepspaces=true,                 
    numbers=left,                    
    numbersep=5pt,                  
    showspaces=false,                
    showstringspaces=false,
    showtabs=false,                  
    tabsize=2
}
\lstset{style=mystyle}

\addbibresource{references.bib}

\begin{document}

\begin{titlepage}
    \centering
    \vspace*{1cm}
    {\Large\textbf{POLITECHNIKA / UCZELNIA TECHNICZNA}\par}
    {\large Wydział Informatyki\par}
    \vspace{3cm}
    {\LARGE\textbf{\MakeUppercase{Tytuł Pracy Dyplomowej}}\par}
    \vspace{2cm}
    {\large\textbf{PRACA DYPLOMOWA INŻYNIERSKA}\par}
    {\large Kierunek: Informatyka Stosowana\par}
    \vspace{3cm}
    \begin{flushleft}
        \textbf{Autor:} Jan Kowalski \\
        \textbf{Promotor:} dr inż. Promotor Akademicki
    \end{flushleft}
    \vfill
    {\large Warszawa, 2027\par}
\end{titlepage}

\tableofcontents
\newpage

\chapter{Wstęp i cel pracy}
Treść wstępu...

\chapter{Przegląd literatury i stan wiedzy (SOTA)}
Treść stanu wiedzy...

\chapter{Projekt architektury systemu IT}
Treść architektury...

\chapter{Implementacja i testy}
Treść implementacji...

\chapter{Badania empiryczne i benchmarki}
Treść badań...

\chapter{Podsumowanie i wnioski}
Wnioski końcowe...

\printbibliography[heading=bibintoc,title={Bibliografia}]

\end{document}

